% review: issue #142, comment 31 --- direct, non-marketing opening
This chapter describes the experimental methodology used to evaluate the
parallel Aho--Corasick scanning strategies introduced in the previous chapter.
The evaluation targets three properties of each strategy --- search
performance, scalability with the number of threads, and correctness with
respect to the sequential baseline --- under a shared-memory multi-core
environment and over pattern dictionaries and text corpora representative of
signature-based intrusion detection. Because the workloads of interest produce
automata whose footprint far exceeds the last-level cache, the reported figures
are expressed as relative quantities (speedup, throughput, and parallel
efficiency) rather than as absolute production capacity, so that the
conclusions concern scaling behaviour instead of the peak rating of a single
machine.

% review: issue #143, comment 32 --- chapter roadmap instead of "principles"
The chapter is organized in the order in which the reader needs the
information. It first establishes the experimental environment and software
stack; then characterizes the pattern dictionaries and text corpora and the
role each one plays; states the three-phase execution model and the invariants
that make the parallelization sound; enumerates the evaluated strategies and
the axes along which they compose; separates the performance metrics from the
supporting diagnostic instrumentation; defines the per-phase measurement
protocol; describes how correctness was established before any performance was
reported; and closes with the experimental design, stating each hypothesis, the
experiment that tests it, and the section of the next chapter that reports the
outcome.

\subsecao{Experimental Environment}

% review: issue #144, comment 33 --- separate the three platform roles
The primary measurement platform, and the source of every headline number,
table, and figure in this work, is a single Intel Core i5-1235U processor; the
canonical data set is the formal sweep executed on it on 2026-05-29. The
i5-1235U is a hybrid (heterogeneous) design of the Alder Lake generation: it
integrates two Performance cores (P-cores) with simultaneous multithreading and
eight Efficiency cores (E-cores) without it, exposing twelve logical hardware
threads in total. The two core types differ in maximum clock frequency
(approximately $4.4$~GHz for P-cores and $3.3$~GHz for E-cores), which is the
source of the load-balancing problem addressed by the frequency-weighted
partitioning strategy. The memory hierarchy comprises a private L1 data cache
per core ($48$~KiB on P-cores, $32$~KiB on E-cores), $1.25$~MiB of L2 per P-core
(and $2$~MiB per E-core cluster), and a single $12$~MiB L3 cache shared by all
cores. Main memory is $31$~GiB of dual-channel DDR4, and the aggregate memory
bandwidth --- on the order of a few tens of GB/s for this configuration --- is
the shared resource for which the parallel scans ultimately contend.

The software stack is a purpose-built laboratory written in C11, compiled with
GCC~13.3.0 using \texttt{-O3 -march=native} and an aggressive warning profile,
and linked against the POSIX Threads library \cite{Nichols1996}. Experiments ran
under Linux (kernel~6.17.0). To minimize scheduling and frequency-scaling
artifacts, the CPU frequency governor was set to \texttt{performance} for the
canonical measurement campaign; the topology-aware searchers
(\texttt{pthread\_chunked\_v3} and \texttt{pthread\_chunked\_v3\_flat})
additionally pin each worker thread to a dedicated logical CPU via
\texttt{pthread\_setaffinity\_np}, while the remaining searchers rely on OS
scheduling without explicit binding. The build was additionally validated under
sanitizer instrumentation, namely the Address and Undefined Behavior sanitizers
and the Thread Sanitizer, as discussed in Section~\ref{subsec:correctness}.

Two further runs play strictly secondary roles and do not displace the
canonical i5 sweep. The first is a portability run on an AMD Ryzen 9 9950X
(Zen 5), a homogeneous design with sixteen identical cores and thirty-two
threads (two-way simultaneous multithreading) and a non-unified $64$~MiB L3
organized as two $32$~MiB slices, one per core complex, backed by DDR5; its
sweep, dated 2026-06-25, is used only to assess the external validity of the i5
findings and is reported as a compact portability comparison in the next
chapter, never as a substitute for the canonical numbers. The second is a
reproducibility re-run of the i5 sweep on 2026-06-28, executed headless on a
cooler machine; it is used to characterize run-to-run variance and is likewise
not a source of headline figures. The single-thread baseline reproduces to
within about $3\%$ between the two i5 runs, but the saturated twelve-thread peak
speedups vary substantially across separate invocations --- a limitation made
explicit in the measurement protocol below.

\subsecao{Pattern Dictionaries and Text Corpora}

The workloads were chosen to sweep the automaton footprint across the cache
hierarchy while keeping the input text realistic for the intrusion-detection
domain. Table~\ref{tab:datasets} summarizes the pattern dictionaries, which were
derived from publicly available Snort and Emerging Threats rule sets, and
Table~\ref{tab:corpora} summarizes the text corpora.

% review: issue #145, comment 34; issue #115, comment 35 --- replace the
% platform-dependent "cache regime" column with a platform-independent role
\begin{table}[H]
    \centering
    \caption{Pattern dictionaries, the resulting automaton footprint, and their experimental role.}
    \label{tab:datasets}
    \small
    \begin{tabular}{@{} l r r r p{4.2cm} @{}}
        \toprule
        \textbf{Dictionary} & \textbf{Patterns} & \textbf{States} & \textbf{Footprint} & \textbf{Experimental role} \\
        \midrule
        Snort-100   & 100      & 1{,}939   & 1.93~MiB   & Smallest automaton; lower-bound contrast regime. \\ \addlinespace
        Snort-1k    & $\sim$1{,}000 & 11{,}974 & 11.92~MiB & Intermediate footprint; brackets the throughput cross-over. \\ \addlinespace
        Snort (full) & 4{,}188  & 55{,}479  & 55.23~MiB  & Primary target dictionary (moderate footprint). \\ \addlinespace
        Emerging Threats & 44{,}678 & 508{,}896 & 506.78~MiB & Largest footprint; most aggressive memory stressor. \\
        \bottomrule
    \end{tabular}

    {\footnotesize\raggedright \textbf{Source:} Prepared by the author from Snort and Emerging Threats rule sets.\par}
\end{table}

\begin{table}[H]
    \centering
    \caption{Text corpora used as scanning input.}
    \label{tab:corpora}
    \small
    \begin{tabular}{@{} l r p{7.2cm} @{}}
        \toprule
        \textbf{Corpus} & \textbf{Size} & \textbf{Characterization and role} \\
        \midrule
        SimpleWiki & $\sim$1.19~GiB & Natural-language encyclopedic text with low match density; used as the contrast corpus in the cross-corpus experiment, which isolates the effect of the text content. \\ \addlinespace
        Enron      & $\sim$1.32~GiB & Real e-mail messages with medium match density; the canonical benchmark corpus. \\ \addlinespace
        Enron ($8\times$) & $\sim$10.59~GiB & The Enron corpus replicated eightfold; amortizes the build phase and reduces scheduler noise in the scaling curves. \\
        \bottomrule
    \end{tabular}

    {\footnotesize\raggedright \textbf{Source:} Prepared by the author. The Enron corpus derives from the public Enron Email Dataset.\par}
\end{table}

% review: issue #145, comment 34 --- footprints are platform-independent;
% cache interpretation deferred to the results chapter and tied to a platform
The two smaller dictionaries (Snort-100 and Snort-1k) serve as
\textit{contrast regimes} that bracket the point at which the automaton stops
fitting comfortably in cache; the full Snort dictionary and the Emerging Threats
subset are the actual targets of the study, the latter being the most aggressive
stressor, with an automaton far larger than the last-level cache of any platform
used here. The footprints in Table~\ref{tab:datasets} are reported as absolute,
platform-independent sizes; how each footprint relates to a particular cache
level depends on the machine and is therefore interpreted in the next chapter
against the cache hierarchy of each platform (for the reference i5, the
$12$~MiB shared L3).

\subsecao{Execution Model and Invariants}

Every execution of the laboratory follows three strictly separated phases; this
separation is the foundation of the lock-free design.

\begin{enumerate}
    \item \textbf{Construction.} The patterns are inserted into a trie, and a
    breadth-first (BFS) traversal computes the goto, failure, and
    dictionary-suffix links, together with the precomputed output structures.
    The resulting automaton is then declared immutable. % review: issue #148, comment 38
    \item \textbf{Search.} Each worker thread reads the shared, read-only
    automaton and its assigned portion of the input text, emitting matches into
    a private, thread-local list. No mutex or atomic operation is used on the
    hot path, because the automaton is never modified during this phase.
    \item \textbf{Merge.} After all workers join, the master thread concatenates
    the thread-local lists in a deterministic order. This phase is always
    sequential and accounts for a negligible fraction of the total time, so it is
    not a target for optimization.
\end{enumerate}

This study respects a set of non-negotiable invariants that make the
parallelization sound. The automaton is strictly read-only during the search;
the hot loop uses neither locks nor atomics; matches remain thread-local until
the join; and the partitioning preserves match ownership without duplication.
For any text-level partitioning, the minimum overlap between adjacent segments
is $L_{\max}-1$ bytes, where $L_{\max}$ is the length of the longest pattern --- as established in the previous chapter, the smallest margin that guarantees that a boundary-straddling pattern is detected by exactly one worker. Match ownership is made disjoint by assigning each
boundary-crossing match to a single segment according to its end position.

\subsecao{Evaluated Strategies}

% review: issue #116, comment 40 --- show which axes are alternatives vs composable
The strategies are organized as a set of orthogonal design axes so that each
performance effect can be attributed to a single cause. Every variant is a
combination of one choice per axis, all selectable at run time within the same
binary and all sharing the same automaton. The axes, and which choices are
mutually exclusive versus composable, are summarized in
Table~\ref{tab:variants}.

\begin{table}[H]
    \centering
    \caption{Evaluated strategies as orthogonal design axes; the work-distribution
    options are alternatives, while the emission and construction axes compose with any of them.}
    \label{tab:variants}
    \small
    \begin{tabular}{@{} l p{9cm} @{}}
        \toprule
        \textbf{Axis} & \textbf{Options (one per run; mutually exclusive within the axis)} \\
        \midrule
        Work distribution & Sequential baseline; static homogeneous chunking; topology-aware weighted chunking; dynamic chunk dispatch (bag of tasks); dictionary sharding by pattern prefix; two-dimensional sharding $\times$ chunking. \\ \addlinespace
        Emission layout & Chained (pointer-chasing) \emph{or} flattened contiguous output table; composes with any work-distribution choice. \\ \addlinespace
        Construction & Sequential build \emph{or} level-synchronous parallel build (BFS); composes with any work-distribution choice. \\
        \bottomrule
    \end{tabular}

    {\footnotesize\raggedright \textbf{Source:} Prepared by the author.\par}
\end{table}

Two of the three axes compose freely. The emission layout (chained
pointer-chasing versus a flattened, contiguous output table) and the
construction mode (sequential versus level-synchronous parallel build) can be
paired with any work-distribution choice and with the sequential baseline,
because the parallel build produces an automaton bit-for-bit identical to the
sequential one and the emission layout does not change which matches are
produced. The work-distribution axis, in contrast, is mutually exclusive: a run
partitions the work in exactly one way. The named searchers in the laboratory
are therefore compositions of these axes --- for example, static chunking with
the flat layout, or topology-aware chunking with the flat layout --- which is
precisely what allows one axis to be varied while the others are held fixed.

% review: issue #149, comment 39 --- topology-aware vs dynamic are distinct axes
The work-distribution options differ in how they assign work to threads. Static
homogeneous chunking splits the text into equal segments, one per thread, with
the boundary overlap required for correctness. Topology-aware chunking and
dynamic dispatch are two \textit{distinct} answers to the load imbalance that
equal segments cause on a heterogeneous processor, and they are not the same
mechanism. Topology-aware chunking sizes each segment \textit{statically}, in
proportion to the maximum clock frequency of the core that will execute it ---
read once from the operating system's per-core frequency information --- and
pins the worker to that core, so that a faster P-core receives a proportionally
larger segment before the scan begins. Dynamic dispatch instead splits the text
into many small, uniform tasks that idle workers pull from a shared atomic
counter \textit{at run time} (a bag-of-tasks scheme), correcting imbalance
reactively rather than from a hardware model. The dictionary-parallelism options
partition the pattern set instead of the text: prefix sharding builds $K$
independent sub-automata, and the two-dimensional composition combines sharding
with text chunking to test whether the two kinds of parallelism add. An
exploratory software-prefetch variant of the chunked searcher also exists in the
laboratory as a microarchitectural probe; it introduces no new axis and is
excluded from the controlled grid and from the reported results.

\subsecao{Metrics and Instrumentation}

% review: issue #117, comment 41 --- separate performance metrics from diagnostics
The performance metrics follow the standard definitions introduced in the
theoretical framework. Search throughput is reported in MB/s (and, where useful,
in Gbps); speedup is the ratio of sequential to parallel search time, as in
Eq.~\ref{eq:speedup}; and parallel efficiency is the speedup normalized by the
thread count, as in Eq.~\ref{eq:efficiency}. Search time is the quantity these
are derived from, and it is always measured separately from construction time,
so that the cost of rebuilding the automaton --- relevant to rule reloading ---
is never conflated with scanning throughput.

A second group of quantities does not measure performance but supports its
interpretation, and is reported as diagnostic instrumentation rather than as a
result. Per-thread timings and the coefficient of variation across repetitions
measure load balance and stability, respectively: the per-thread spread exposes
synchronization waste at the final barrier, while the coefficient of variation
quantifies how noisy a measurement is rather than how fast it is. Match counts
are recorded for every run as a correctness check, not a performance figure (see
Section~\ref{subsec:correctness}). Construction (build) time is recorded for the
parallel-build experiment. Finally, the laboratory writes every run, together
with a snapshot of its environment, to a structured log that is ingested into an
SQLite database, which is the substrate for all aggregation and cross-checking;
the automaton footprint in bytes is retained there too --- not as a headline
figure, but as the causal variable used to explain throughput across regimes.

\subsecao{Measurement Protocol}

Each reported data point follows the exact per-phase repetition protocol
summarized in Table~\ref{tab:protocol}. 

\begin{table}[H]
    \centering
    \caption{Per-phase repetition protocol of the formal sweep (2026-05-29), verified against \texttt{sweep.db}.}
    \label{tab:protocol}
    \small
    \begin{tabular}{@{} c l r r @{}}
        \toprule
        \textbf{Phase} & \textbf{Description} & \textbf{Warm-ups} & \textbf{Timed iters} \\
        \midrule
        A & Speedup curves        & 2 & 5 \\
        B & Footprint sweep       & 2 & 5 \\
        C & Cross-corpus          & 2 & 5 \\
        D & Per-thread diagnostic & 1 & 3 \\
        E & Parallel construction & 0 & 1 \\
        \bottomrule
    \end{tabular}

    {\footnotesize\raggedright \textbf{Source:} Prepared by the author from \texttt{sweep.db} (sweep 2026-05-29).\par}
\end{table}

Phase~A (speedup curves) sweeps thread
counts $T \in \{1,2,4,6,8,10,12\}$ for the canonical Enron corpus and
$T \in \{1,4,8,12\}$ for the eight-fold replica, as recorded in
\texttt{sweep.db}. Phases A--C (speedup curves, footprint, and cross-corpus)
each execute two warm-up iterations followed by five timed iterations; phase D (per-thread diagnostic) executes one warm-up followed by
three timed iterations; and phase E (parallel construction) executes a single
timed build with no warm-up. The warm-up iterations pay for page faults and warm
the caches so that the timed iterations measure steady-state behaviour. Both the
mean and the best (minimum) times are recorded; the mean is used for the reported
throughput, and the dispersion between iterations is retained as an indicator of
stability.

A caveat specific to the measurement platform must be stated explicitly. The
reference machine is a thermally constrained mobile processor.
Under sustained all-core load lasting hours, the package throttles its
frequency, and the measured multi-threaded throughput settles roughly $15$ to
$35\%$ below what isolated bursts on a cold processor report. The values
published in this work are deliberately the sustained, under-load figures, which
are conservative and reproducible; cold-burst peaks are treated as the
non-sustainable upper bound and are not used as headline results. This thermal
sensitivity is itself a finding consistent with the central thesis: at the scale
of interest, performance is governed by physical resource limits rather than by
the number of threads.

\subsecao{Correctness Validation}
\label{subsec:correctness}

Performance is only meaningful once correctness is established, so every parallel
variant was required to reproduce the output of the sequential baseline before
being benchmarked. Correctness was checked at three complementary levels. First,
the automated test suite requires every strategy to reproduce the exact match
\emph{set} of the sequential baseline---each match identified by its end position
and pattern identifier---across thread counts of $\{1,2,3,4,7,8\}$; the same suite
also verifies that the level-synchronous parallel construction produces an
automaton bit-for-bit identical to the sequentially built one, so that swapping in
the parallel build can never alter the reported matches. Second, on the full-scale
workloads, the MD5 digest of the canonicalized match stream---the full set of
matches placed in a deterministic order by end position and pattern
identifier---was computed for each parallel strategy and compared against the
digest of the sequential baseline. An agreement certifies the exact \emph{multiset}
of matches (their count, their end positions, and their pattern identifiers), a far
stronger guarantee than match-count parity alone; it was confirmed for every
strategy up to a thread count of $64$ (a $5.3\times$ oversubscription of the
hardware), demonstrating that the partitioning logic is robust under heavy
oversubscription. For the static-chunking strategies, whose thread-local results
are merged in text order, the digest of the raw emission stream agrees as well; the
dynamic-dispatch and dictionary-sharded strategies emit the identical multiset in a
different order, so for them only the canonicalized digest coincides. Finally,
ThreadSanitizer reported no data races during the parallel search, empirically
confirming the read-only-automaton invariant and the absence of unsynchronized
shared writes on the hot path.

\subsecao{Experimental Design and Hypotheses}

The experiments were structured to test a set of explicit hypotheses, several of
which anticipate negative or qualified outcomes. The main hypotheses are: that
text chunking scales sublinearly on the hybrid processor but suffers from
load imbalance when segments are sized uniformly; that the automaton footprint,
rather than the corpus content, governs throughput once the working set escapes
the cache; that the flattened output layout reduces emission cost in both the
sequential and parallel cases; that a level-synchronous parallel construction
reduces build time for large dictionaries; that prefix-based dictionary sharding
can recover throughput in the memory-bound regime by shrinking the per-shard
working set; and that a two-dimensional composition of sharding and chunking
could combine cache locality with text parallelism. Each hypothesis maps to a
controlled experiment that varies a single axis while holding the others fixed,
always comparing against the sequential baseline for headline speedup and against
the strongest text-parallel strategy when the question concerns composition.
Reporting the strategies that fail to outperform their simpler counterparts is
treated as a scientific result in its own right, since it delimits the boundary
beyond which plausible optimizations stop paying off on this class of hardware.
